% Part: first-order-logic
% Chapter: completeness
% Section: outline

\documentclass[../../../include/open-logic-section]{subfiles}

\begin{document}

\iftag{FOL}
      {\olfileid{fol}{com}{out}}
      {\olfileid{pl}{com}{out}}

\olsection{مخطط البرهان}

برهان مبرهنة الاكتمال معقد بعض الشيء، ومن السهل أن يتيه القارئ عند
قراءته أول مرة. فلنضع إذن مخططًا للبرهان. والخطوة الأولى هي تغيير
وجهة النظر بحيث ينكشف لنا طريق إلى البرهان. إذا فكرنا في الاكتمال على
صورة «كلما كان $\Gamma \Entails !A$ كان $\Gamma \Proves !A$»، فقد
يصعب حتى أن تخطر لنا فكرة: فلكي نبين أن $\Gamma \Proves !A$ علينا أن
نجد !!a{derivation}، ولا يبدو أن الفرض $\Gamma \Entails !A$ يعيننا على
ذلك بوجه. ويمكن في بعض الأنساق البرهانية إنشاء !!a{derivation} مباشرة،
لكننا سنتبع نهجًا مختلفًا قليلًا. وتغيير المنظور المطلوب هو الآتي:
يمكن أيضًا صياغة الاكتمال هكذا: «إذا كانت $\Gamma$ متسقة فهي قابلة
للإرضاء». فلعلنا نستطيع أن نستعمل المعلومات في~$\Gamma$ مع فرض اتساقها
لننشئ \iftag{FOL}{!!a{structure}}{!!a{valuation}} ترضي كل
\iftag{FOL}{!!{sentence}}{!!{formula}} في~$\Gamma$. فنحن نعرف في نهاية
الأمر نوع ال\iftag{FOL}{!!{structure}}{!!{valuation}} التي نبحث عنها:
واحدة تكون كما تصفها $\Gamma$!{}

إذا لم تحتو $\Gamma$ إلا على \iftag{FOL}{!!{sentence}s
  ذرية}{!!{propositional variable}s}، سهل إنشاء نموذج لها.
\iftag{FOL}{ افترض أن كل ال!!{sentence}s الذرية من الصورة
  $\Atom{P}{a_1,\dots,a_n}$، حيث تكون $a_i$ !!{constant}s.}{} وكل ما
علينا فعله هو إيجاد
\iftag{FOL}{%
  !!a{domain}~$\Domain{M}$ وإسناد لـ~$P$ بحيث
  $\Sat{M}{\Atom{P}{a_1,\ldots,a_n}}$. وليس ذلك صعبًا: ضع
  $\Domain{M} = \Nat$ و$\Assign{\Obj c_i}{M} = i$، ولكل
  $\Atom{P}{a_1,\ldots,a_n} \in \Gamma$ ضع الصف $\tuple{k_1,
    \dots, k_n}$ في $\Assign{P}{M}$، حيث $k_i$ فهرس الرمز
  الثابت~$a_i$ (أي إن $a_i \ident \Obj c_{k_i}$).}{%
  !!a{valuation}~$\pAssign{v}$ بحيث $\pSat{v}{p}$ لكل $p \in
  \Gamma$. ولنفعل ذلك بوضع $\pAssign{v}(p) = \True$ إذا وفقط إذا كان
  $p \in \Gamma$.}

افترض الآن أن $\Gamma$ تحتوي !!{formula} ما~$\lnot !B$، حيث $!B$
ذرية. وقد نقلق من أن يتعارض إنشاء
$\iftag{FOL}{\Struct{M}}{\pAssign{v}}$ مع إمكان جعل $\lnot !B$ صادقة.
وهنا تظهر فائدة اتساق~$\Gamma$: إذا كان $\lnot !B \in \Gamma$، فإن
$!B \notin \Gamma$، وإلا كانت $\Gamma$ غير متسقة. وإذا كان
$!B \notin \Gamma$، فبحسب إنشائنا لـ
$\iftag{FOL}{\Struct{M}}{\pAssign{v}}$ لدينا
$\iftag{FOL}{\Sat/{M}}{\pSat/{v}}{!B}$، ومن ثم
$\iftag{FOL}{\Sat{M}}{\pSat{v}}{\lnot !B}$. وهذا جيد حتى الآن.

ولكن ماذا لو احتوت $\Gamma$ صيغًا مركبة غير ذرية؟ لنقل إنها تحتوي
$!A \land !B$. لكي نجعلها صادقة ينبغي أن نسير كما لو كانت كل من $!A$
و$!B$ في~$\Gamma$. وإذا كان $!A \lor !B \in \Gamma$، فعلينا أن نجعل
واحدة منهما على الأقل صادقة؛ أي أن نسير كما لو كانت إحداهما
في~$\Gamma$.

يوحي هذا بالفكرة الآتية: نضيف !!{formula}s أخرى إلى~$\Gamma$ بحيث
(أ)~نبقي المجموعة الناتجة متسقة، و(ب)~نضمن أنه لكل
\iftag{FOL}{!!{sentence}}{!!{formula}} ممكنة~$!A$، إما أن تكون $!A$
في المجموعة الناتجة أو أن تكون
$\lnot !A$ فيها، و(ج)~نضمن أنه كلما كانت $!A \land !B$ في المجموعة
كانت كل من $!A$ و$!B$ فيها، وإذا كانت $!A \lor !B$ فيها كانت واحدة
على الأقل من $!A$ و$!B$ فيها أيضًا، وهكذا. ونواصل فعل ذلك (إلى الأبد
ربما). ولنسم مجموعة كل ال!!{formula}s المضافة على هذا النحو~$\Gamma^*$.
عندئذ يزودنا إنشاؤنا أعلاه بــ
\iftag{FOL}{!!a{structure}~$\Struct{M}$}{!!a{valuation}~$\pAssign{v}$}
يمكننا أن نبرهن بالاستقراء أنها ترضي كل
\iftag{FOL}{ال!!{sentence}s}{ال!!{formula}s} في~$\Gamma^*$، ومن ثم كل
\iftag{FOL}{ال!!{sentence}s}{ال!!{formula}s} في~$\Gamma$ أيضًا، لأن
$\Gamma \subseteq \Gamma^*$. ويتبين أن ضمان (أ) و(ب) كافٍ. وتسمى
\iftag{FOL}{مجموعة ال!!{sentence}s}{مجموعة ال!!{formula}s} التي يتحقق
فيها (ب) \emph{كاملة}. ومن ثم ستكون مهمتنا توسيع المجموعة المتسقة
~$\Gamma$ إلى مجموعة متسقة وكاملة~$\Gamma^*$.

\iftag{FOL}{%
في هذه الخطة عقبة واحدة: إذا كان $\lexists[x][!A(x)] \in \Gamma$،
فسنأمل أن نستطيع اختيار !!{constant} ما~$c$ وإضافة $!A(c)$ أثناء هذه
العملية. ولكن كيف نعلم أن ذلك ممكن دائمًا؟ فلعل لغتنا لا تحتوي إلا
قليلًا من ال!!{constant}s، ولكل واحد منها $\lnot !A(c) \in \Gamma$.
ولا يمكننا أن نضيف $!A(c)$ أيضًا، لأن ذلك يجعل المجموعة غير متسقة،
ولن نعرف حينئذ هل يجب على $\Struct{M}$ أن تجعل $!A(c)$ أم
$\lnot !A(c)$ صادقة. بل قد يحدث أن لا تحتوي $\Gamma$ إلا جملًا في لغة
لا رموز ثابتة فيها أصلًا (كلغة نظرية المجموعات مثلًا).

وحل هذه المشكلة هو ببساطة أن نضيف في البداية عددًا لا نهائيًا من
الرموز الثابتة، ومعها جملًا تصل بينها وبين الكمّيات على الوجه الصحيح.
(وعلينا طبعًا أن نتحقق من أن ذلك لا يمكن أن يُدخل تناقضًا.)

يعمل إنشاؤنا الأصلي جيدًا إذا لم يكن في الجمل الذرية إلا
!!{constant}s. لكن قد تحتوي اللغة أيضًا !!{function}s. وقد يصعب عندئذ
العثور على الدوال المناسبة على~$\Nat$ لإسنادها إلى هذه ال!!{function}s
بحيث يعمل كل شيء. وهذه حيلة أخرى: بدلًا من استعمال $i$ لتفسير
$\Obj c_i$، اتخذ مجموعة ال!!{constant}s نفسها مجالًا. وعندئذ تستطيع
$\Struct M$ أن تسند كل !!{constant} إلى نفسه:
$\Assign{\Obj c_i}{M} = \Obj c_i$. بل لم لا نمضي إلى النهاية ونجعل
$\Domain{M}$ مجموعة جميع \emph{حدود} اللغة!{} إذا فعلنا ذلك، وُجد إسناد
واضح للدوال (التي تأخذ حدودًا حججًا وتكون قيمها حدودًا) إلى
ال!!{function}s: نسند إلى ال!!{function}~$\Obj f^n_i$ الدالة التي تأخذ
$n$ من الحدود $t_1$، \dots،~$t_n$ مدخلات، وتنتج الحد
$\Obj f^n_i(t_1, \dots, t_n)$ قيمةً.

وآخر قطعة في الأحجية هي ما ينبغي فعله بـ~$\eq$. لل!!{predicate}~$\eq$
تفسير ثابت: يتحقق $\Sat{M}{\eq[t][t']}$ إذا وفقط إذا كان
$\Value{t}{M} = \Value{t'}{M}$. فإذا رتبنا الأمور بحيث تكون !!{value}
الحد~$t$ هي $t$ نفسه، فلن تجعل هذه ال!!{structure} \emph{أي} جملة من
الصورة $\eq[t][t']$ صادقة إلا إذا كان $t$ و$t'$ الحد نفسه. وهذه مشكلة
طبعًا، لأن كل نظرية مهمة تقريبًا في لغة ذات !!{function}s تتخذ
$\eq[t][t']$ مبرهنةً مع أن $t$ و$t'$ ليسا الحد نفسه (كما في نظريات
الحساب: $\eq[(\Obj 0+ \Obj 0)][\Obj 0]$). ولحل هذه المشكلة نغير مجال
$\Struct M$: فبدلًا من استعمال الحدود أشياءً في~$\Domain{M}$، نستعمل
مجموعات من الحدود، وتكون كل مجموعة بحيث تحتوي جميع الحدود التي تتطلب
جمل~$\Gamma$ تساويها. فإذا كانت $\Gamma$ مثلًا نظرية في الحساب، فستحتوي
إحدى هذه المجموعات: $\Obj 0$ و$(\Obj 0 + \Obj 0)$ و$(\Obj
0 \times \Obj 0)$، وهكذا. وهذه هي المجموعة التي نسندها إلى $\Obj 0$،
وسيتبين أنها أيضًا قيمة كل الحدود التي فيها، ومنها مثلًا
$(\Obj 0 + \Obj 0)$. ولذلك تكون الجملة
$\eq[(\Obj 0+ \Obj 0)][\Obj 0]$ صادقة في هذه ال!!{structure} المعدّلة.}{}

وإليك الآن ما سنفعله. نبحث أولًا خصائص المجموعات المتسقة
ال!!{complete}، ونبرهن بوجه خاص أن المجموعة المتسقة !!a{complete}
تحتوي $!A \land !B$ إذا وفقط إذا كانت تحتوي كلًا من $!A$ و~$!B$،
وتحتوي $!A \lor !B$ إذا وفقط إذا كانت تحتوي واحدة منهما على الأقل،
وهكذا (\olref[ccs]{prop:ccs}).\iftag{FOL}{ ثم نعرّف مجموعات الجمل
  «المشبعة» ونبحثها. والمجموعة المشبعة هي التي تحتوي شرطيات تصل كل
  !!{sentence} مكمّمة بحالاتها (\olref[hen]{defn:henkin-exp}). ونبين أن
  كل مجموعة متسقة~$\Gamma$ يمكن دائمًا توسيعها إلى مجموعة مشبعة
  ~$\Gamma'$ (\olref[hen]{lem:henkin}). وإذا كانت مجموعة ما متسقة
  ومشبعة و!!{complete}، فلها أيضًا خاصية أنها تحتوي
  \iftag{prvEx}{$\lexists[x][!A(x)]$ إذا وفقط إذا كانت تحتوي $!A(t)$
    لبعض الحدود المغلقة~$t$}{}\iftag{defEx,defAll}{}{، و
  }\iftag{prvAll}{$\lforall[x][!A(x)]$ إذا وفقط إذا كانت تحتوي $!A(t)$
    لكل الحدود المغلقة~$t$}{} (\olref[hen]{prop:saturated-instances}).}{}
ثم نأخذ المجموعة المتسقة\iftag{FOL}{ المشبعة}{}
~$\iftag{FOL}{\Gamma'}{\Gamma}$، ونبين إمكان توسيعها إلى مجموعة
\iftag{FOL}{مشبعة ومتسقة و}{متسقة و}!!{complete}~$\Gamma^*$
(\olref[lin]{lem:lindenbaum}). وسنستعمل هذه المجموعة $\Gamma^*$ لتعريف
\iftag{FOL}{نموذج الحدود~$\Struct M(\Gamma^*)$}{!!{valuation}
~$\pAssign v(\Gamma^*)$}. \iftag{FOL}{مجال نموذج الحدود هو مجموعة
  الحدود المغلقة، ويحدد ال!!{sentence}s الذرية تفسير !!{predicate}sه}{يحدد
  ال!!{propositional variable}s التقييم} في~$\Gamma^*$
(\olref[mod]{defn:termmodel}). وسنستعمل خصائص المجموعات
\iftag{FOL}{المشبعة و}{}المتسقة الكاملة لنبيّن أنه بالفعل
$\iftag{FOL}{\Sat{M(\Gamma^*)}}{\pSat{v(\Gamma^*)}}{!A}$ إذا وفقط إذا
كان $!A \in \Gamma^*$ (\olref[mod]{lem:truth})، ولذلك بوجه خاص
$\iftag{FOL}{\Sat{M(\Gamma^*)}}{\pSat{v(\Gamma^*)}}{\Gamma}$.
\iftag{FOL}{وأخيرًا سننظر في كيفية تعريف نموذج حدود حين تحتوي
  $\Gamma$ الرمز~$\eq$ أيضًا (\olref[ide]{defn:term-model-factor})،
  ونبين أنه يرضي~$\Gamma^*$ (\olref[ide]{lem:truth}).}{}

\end{document}
